\centerline{\bf \Huge Сириус}


\begin{flushright}
        \scriptsize{}
\end{flushright}


\problem{1}
В строку записаны 5 чисел. Могло ли оказаться так, что сумма любых двух соседних чисел - отрицательна, а сумма всех пяти чисел - положительна?

\problem{2}
Найдите 5 последних цифр числа $3+33+333+\ldots+3 \ldots 3$ (в последнем слагаемом 2019 троек).

\problem{3}
B треугольнике ABC , в котором AC в два раза больше AB , провели биссектрису AK. Оказалось, что AK=KC. Найти угол ABC.


\problem{4}
На столе лежат 10 карточек с числами: $1 ;-2 ; 3 ;-4 ; 5 ;-6 ; 7 ;-8 ; 9 ;-10$. Двое по очереди выбирают себе по одной карточке, и, после того, как карточки закончатся, подсчитывают суммы чисел на своих карточках. Выигрывает игрок, у которого получилась сумма, большая по модулю. Сможет ли начинающий выиграть в этой игре независимо от ходов второго?

\problem{5}
Существует ли натуральное составное число n , большее 1000 , такое, что любой простой делитель числа n больше, чем любой простой делитель числа $\mathrm{n}+1$ ?

\problem{6}
У Васи и Пети есть по $n$ монет. Каждая монета имеет достоинство $1$ копейка, $50$ копеек или $1$ рубль. Оказалось, что у мальчиков поровну денег, но наборы монет не совпадают. Не обязательно, что монеты каждого достоинства присутствуют. При каком наименьшем n такое возможно?